% 第4章：不确定性与概率规划
\chapter{不确定性与概率规划}

\begin{levelbox}
本章介绍在不确定环境下的规划方法，包括条件规划、MDP/POMDP和强化学习基础。通过案例v3版本学习如何处理观测和状态不确定性。建议学时：8学时。
\end{levelbox}

% =============================================
\section{\caseA{v3}云层遮挡下的观测规划}
% =============================================

\begin{caseboxA}[案例A-v3：不确定性观测环境]
\textbf{场景演进}：v2假设观测必定成功，v3引入不确定性。

\textbf{新增不确定因素}：
\begin{itemize}
  \item \textbf{云层遮挡}：观测成功与否取决于云量，概率可预测但不确定
  \item \textbf{设备故障}：成像设备有一定故障概率
  \item \textbf{目标移动}：部分目标可能在观测窗口内移动
\end{itemize}

\textbf{为什么v2方法不够？}
\begin{itemize}
  \item 确定性规划假设动作效果完全可预测
  \item 无法处理"观测失败后重新规划"的场景
  \item 需要概率建模和期望收益优化
\end{itemize}

\textbf{预习思考}：
\begin{enumerate}
  \item 如何建模"观测可能失败"这一不确定性？
  \item 是否应该为每个目标安排备选观测窗口？
  \item 如何在不确定性下最大化期望观测收益？
\end{enumerate}
\end{caseboxA}

% =============================================
\section{条件规划与一致性规划}
% =============================================

\subsection{非确定性动作}

\begin{definition}[非确定性动作]
非确定性动作的效果不唯一，执行后可能到达多个可能状态之一：
\[ \text{Effects}(a) = \{e_1, e_2, \ldots, e_k\} \]
具体到达哪个状态在执行前未知。
\end{definition}

\subsection{条件规划}

\textbf{条件规划}（Contingent Planning）允许在执行过程中进行观测和分支。

\begin{keypoint}[条件规划的输出]
条件规划的输出不是线性动作序列，而是一棵\textbf{条件规划树}：
\begin{itemize}
  \item 内部节点：感知动作，根据观测结果分支
  \item 叶节点：目标状态或失败
  \item 执行时根据实际观测选择分支
\end{itemize}
\end{keypoint}

\begin{example}[案例A的条件规划]
\begin{verbatim}
1. 尝试观测目标T1
2. 感知：观测是否成功？
   - 如果成功：继续观测T2
   - 如果失败（云层遮挡）：
     - 等待下一观测窗口
     - 或转向备选目标T1'
\end{verbatim}
\end{example}

\subsection{一致性规划}

\textbf{一致性规划}（Conformant Planning）：在不确定初始状态下，寻找一个在所有可能初始状态下都能达成目标的动作序列。

\begin{itemize}[leftmargin=2em]
  \item 不依赖执行时观测
  \item 在信念空间（Belief Space）中搜索
  \item 计算复杂度高（EXPSPACE-完全）
\end{itemize}

\subsection{HTN规划中的不确定性}

值得注意的是，不确定性同样存在于层次任务网络规划中。华中科技大学王红卫教授、祁超教授团队将不确定HTN规划中的不确定性归纳为三类：\textbf{状态不确定}、\textbf{动作效果不确定}和\textbf{任务分解不确定}（参见Liu等人在Knowledge-Based Systems 2016发表的论文）。这一分类框架将第3章的HTN方法与本章的概率规划方法有机关联——HTN的层次分解结构可与概率推理相结合，为应急响应等不确定性环境下的复杂任务提供更具表达力的规划模型。

% =============================================
\section{MDP/POMDP框架}
% =============================================

\subsection{马尔可夫决策过程（MDP）}

\begin{definition}[MDP]
MDP由五元组 $\langle S, A, T, R, \gamma \rangle$ 定义：
\begin{itemize}
  \item $S$：状态空间
  \item $A$：动作空间
  \item $T(s'|s,a)$：状态转移概率
  \item $R(s,a)$：即时奖励函数
  \item $\gamma \in [0,1)$：折扣因子
\end{itemize}
\end{definition}

\textbf{目标}：找到策略 $\pi: S \rightarrow A$，最大化期望累积奖励：
\[ V^\pi(s) = \mathbb{E}\left[\sum_{t=0}^{\infty} \gamma^t R(s_t, \pi(s_t)) \mid s_0 = s\right] \]

\subsection{值迭代与策略迭代}

\begin{algbox}[值迭代算法]
\begin{enumerate}
  \item 初始化 $V(s) = 0$ 对所有 $s$
  \item 重复直到收敛：
  \[ V(s) \leftarrow \max_{a \in A} \left[ R(s,a) + \gamma \sum_{s'} T(s'|s,a) V(s') \right] \]
  \item 提取最优策略：
  \[ \pi^*(s) = \arg\max_{a} \left[ R(s,a) + \gamma \sum_{s'} T(s'|s,a) V(s') \right] \]
\end{enumerate}
\end{algbox}

\subsection{部分可观测MDP（POMDP）}

\begin{definition}[POMDP]
POMDP扩展MDP，增加观测模型：
\begin{itemize}
  \item $\Omega$：观测空间
  \item $O(o|s',a)$：在执行动作 $a$ 到达状态 $s'$ 后观测到 $o$ 的概率
\end{itemize}
\end{definition}

POMDP中智能体不能直接观测状态，而是维护一个\textbf{信念状态}（Belief State）$b$——关于当前状态的概率分布。

\subsection{案例A-v3的POMDP建模}

\begin{example}[云层遮挡的POMDP模型]
\textbf{状态}：$s = (\text{已观测目标集合}, \text{各目标云量状态})$

\textbf{动作}：$a = \text{observe}(t)$ 或 $\text{wait}$

\textbf{转移}：
\begin{itemize}
  \item 云量状态随时间演化（马尔可夫过程）
  \item 观测成功概率取决于云量
\end{itemize}

\textbf{观测}：
\begin{itemize}
  \item 执行observe后观测到"成功"或"失败"
  \item 云量状态不直接可观测，需通过气象预报推断
\end{itemize}

\textbf{奖励}：成功观测目标获得奖励，与目标优先级成正比
\end{example}

% =============================================
\section{\caseB{v3}战场迷雾下的决策}
% =============================================

\begin{caseboxB}[案例B-v3：不完全信息作战]
\textbf{场景演进}：v2假设敌情完全已知，v3引入"战场迷雾"。

\textbf{新增不确定性}：
\begin{itemize}
  \item \textbf{敌方位置不确定}：只知道敌方可能在某些区域
  \item \textbf{敌方意图不明}：不知道敌方是否会机动
  \item \textbf{侦察结果不完全}：侦察可能漏报或误报
\end{itemize}

\textbf{POMDP建模}：
\begin{itemize}
  \item 状态包含敌方真实位置和状态（不可直接观测）
  \item 侦察动作提供关于敌情的噪声观测
  \item 需要在信念空间中规划
\end{itemize}
\end{caseboxB}

% =============================================
\section{强化学习与世界模型}
% =============================================

\subsection{强化学习基础}

当转移概率 $T$ 和奖励函数 $R$ 未知时，可以通过\textbf{强化学习}（RL）从交互经验中学习策略。

\begin{itemize}[leftmargin=2em]
  \item \textbf{基于值的方法}：学习值函数 $Q(s,a)$，如Q-learning、DQN
  \item \textbf{基于策略的方法}：直接学习策略参数，如REINFORCE、PPO
  \item \textbf{Actor-Critic}：同时学习策略和值函数
\end{itemize}

\subsection{基于模型的强化学习}

\begin{keypoint}[基于模型的RL与规划]
基于模型的强化学习先学习环境模型（世界模型），再基于模型进行规划：
\begin{enumerate}
  \item 从经验中学习转移模型 $\hat{T}$ 和奖励模型 $\hat{R}$
  \item 在学习的模型中进行规划（如值迭代、蒙特卡洛树搜索）
  \item 执行规划结果，收集新经验
\end{enumerate}
\end{keypoint}

\subsection{MuZero与Dreamer}

\textbf{MuZero}（DeepMind, 2020）：
\begin{itemize}[leftmargin=2em]
  \item 学习隐状态表示和动态模型
  \item 在隐空间中进行蒙特卡洛树搜索
  \item 不需要环境的完整规则
  \item 在围棋、Atari等任务上达到超人水平
\end{itemize}

\textbf{Dreamer系列}（Hafner et al.）：
\begin{itemize}[leftmargin=2em]
  \item 学习世界模型（编码器+动态模型+解码器）
  \item 在"想象"（世界模型中的模拟）中训练策略
  \item DreamerV3在多种任务上表现优异
\end{itemize}

% =============================================
\section{本章小结与讨论}
% =============================================

\subsection*{本章小结}

\begin{itemize}[leftmargin=2em]
  \item 条件规划通过感知和分支处理不确定性
  \item MDP/POMDP是不确定性规划的数学框架
  \item 强化学习可以在未知环境中学习策略
  \item 世界模型（MuZero、Dreamer）结合了学习和规划
\end{itemize}

\begin{discussbox}
\textbf{讨论题}：
\begin{enumerate}
  \item 案例A-v3中，如果云量预报完全准确，问题是否退化为确定性规划？
  \item POMDP的精确求解是PSPACE-完全的，实际中如何处理大规模POMDP？
  \item 案例B-v3中，侦察误报（将非目标判为目标）和漏报（未发现目标）哪个更危险？如何在建模中体现？
  \item MuZero和传统规划器有什么本质区别？各有什么优缺点？
\end{enumerate}

\textbf{编程练习}：
\begin{enumerate}
  \item 将案例A-v3简化为小规模MDP，使用值迭代求解
  \item 使用POMDP-solve或SARSOP求解案例A-v3的POMDP
  \item 比较MDP策略和条件规划在模拟环境中的表现
\end{enumerate}
\end{discussbox}
